\chapter{Singleton bedingte Unabhängigkeit}
\label{ch:singl-ci}

\section{Problemstellung}
\label{sec:singl-ci-def}
Im Folgenden betrachten wir eine eingeschränkte Variante des Entscheidungsproblems
der \textit{stable}-Unabhängigkeit. Sei \(F=(A,R)\) ein AF, seien
\(x,y\in A\) zwei verschiedene Argumente und sei
\(Z\subseteq A\setminus\{x,y\}\). Damit sind
\(X=\{x\}\), \(Y=\{y\}\) und \(Z\) paarweise disjunkt. Zu entscheiden ist, ob \(\{x\}\bot_{\mathrm{st}}\{y\}\mid Z\)
gilt. Das zugehörige Entscheidungsproblem lautet:
\[
\begin{array}{ll}
\textit{SingInd}_{\mathrm{st}}\\[2mm]
\text{Input:}  & \text{AF } F=(A,R),\ \text{Argumente } x,y\in A
                 \text{ mit } x\neq y,\ \text{und } Z\subseteq A\setminus\{x,y\}.\\[1mm]
\text{Output:} & \text{TRUE} \Longleftrightarrow
                 \{x\}\bot_{\mathrm{st}}\{y\}\mid Z.
\end{array}
\]
Dabei bedeutet \(\{x\}\bot_{\mathrm{st}}\{y\}\mid Z\), dass für alle
\(\lambda_1,\lambda_2\in\Lambda_{\mathrm{st}}(F)\) mit
\(\lambda_1|_Z=\lambda_2|_Z\) ein \(\lambda_3\in\Lambda_{\mathrm{st}}(F)\)
existiert, sodass
    \(\lambda_3(x)=\lambda_1(x),~
    \lambda_3(y)=\lambda_2(y),~
    \lambda_3|_Z=\lambda_1|_Z\)
gilt.

\section{Theoretische Grundlage}
\label{sec:singl-ci-theorie}

Bevor wir den Algorithmus angeben, charakterisieren wir die
\textit{stable}-Unabhängigkeit im Singleton-Fall durch eine einfache
kombinatorische Bedingung. Sie bildet die Grundlage für die Korrektheit des
Algorithmus und rechtfertigt die spätere Beschränkung auf wenige
Kandidatenfälle.

\begin{Definition}
    \label{def:singl-realisierbar}
    Sei \(F=(A,R)\) ein AF, seien \(x,y\in A\) verschieden und sei
    \(Z\subseteq A\setminus\{x,y\}\). Für eine Belegung
    \(\eta:Z\to L_{\mathrm{st}}\) mit
    \(L_{\mathrm{st}}=\{\mathit{in},\mathit{out}\}\) definieren wir die Menge
    der unter \(\eta\) realisierbaren \((x,y)\)-Labelpaare durch
    \[
        R_\eta
        :=
        \bigl\{
        (\lambda(x),\lambda(y))
        \;\big|\;
        \lambda\in\Lambda_{\mathrm{st}}(F),\ \lambda|_Z=\eta
        \bigr\}
        \subseteq L_{\mathrm{st}}^2.
    \]
    Die Belegung \(\eta\) heißt \emph{realisierbar}, falls
    \(R_\eta\neq\emptyset\), falls also ein
    \(\lambda\in\Lambda_{\mathrm{st}}(F)\) mit \(\lambda|_Z=\eta\) existiert.
\end{Definition}

Eine Menge \(S\subseteq L_{\mathrm{st}}^2\) nennen wir \emph{Produktrechteck},
wenn \(P,Q\subseteq L_{\mathrm{st}}\) mit \(S=P\times Q\) existieren.
Gleichwertig ist \(S\) genau dann ein Produktrechteck, wenn aus
\((\alpha_1,\beta_1),(\alpha_2,\beta_2)\in S\) stets
\((\alpha_1,\beta_2)\in S\) folgt.

\begin{Proposition}
    \label{prop:singl-rechteck}
    Sei \(F=(A,R)\) ein AF, seien \(x,y\in A\) verschieden und sei
    \(Z\subseteq A\setminus\{x,y\}\). Dann gilt
    \(\{x\}\bot_{\mathrm{st}}\{y\}\mid Z\) genau dann, wenn \(R_\eta\) für jede
    realisierbare Belegung \(\eta:Z\to L_{\mathrm{st}}\) ein Produktrechteck
    ist.
\end{Proposition}

\begin{proof}
    \label{pr:singl-rechteck}
    ``\(\Rightarrow\)'': Es gelte \(\{x\}\bot_{\mathrm{st}}\{y\}\mid Z\), und
    sei \(\eta\) realisierbar. Seien
    \((\alpha_1,\beta_1),(\alpha_2,\beta_2)\in R_\eta\). Dann existieren
    \(\lambda_1,\lambda_2\in\Lambda_{\mathrm{st}}(F)\) mit
    \(\lambda_1|_Z=\lambda_2|_Z=\eta\) sowie \(\lambda_1(x)=\alpha_1\),
    \(\lambda_1(y)=\beta_1\), \(\lambda_2(x)=\alpha_2\) und
    \(\lambda_2(y)=\beta_2\). Wegen \(\lambda_1|_Z=\lambda_2|_Z\) liefert die
    Unabhängigkeit nach Definition~\ref{def:BU:AFs} ein
    \(\lambda_3\in\Lambda_{\mathrm{st}}(F)\) mit
    \(\lambda_3(x)=\lambda_1(x)=\alpha_1\),
    \(\lambda_3(y)=\lambda_2(y)=\beta_2\) und \(\lambda_3|_Z=\eta\). Also ist
    \((\alpha_1,\beta_2)\in R_\eta\), und \(R_\eta\) ist ein Produktrechteck.

    ``\(\Leftarrow\)'': Seien
    \(\lambda_1,\lambda_2\in\Lambda_{\mathrm{st}}(F)\) mit
    \(\lambda_1|_Z=\lambda_2|_Z=:\eta\). Dann ist \(\eta\) realisierbar, und es
    gilt \((\lambda_1(x),\lambda_1(y)),(\lambda_2(x),\lambda_2(y))\in R_\eta\).
    Da \(R_\eta\) ein Produktrechteck ist, folgt
    \((\lambda_1(x),\lambda_2(y))\in R_\eta\); es existiert also ein
    \(\lambda_3\in\Lambda_{\mathrm{st}}(F)\) mit \(\lambda_3|_Z=\eta\),
    \(\lambda_3(x)=\lambda_1(x)\) und \(\lambda_3(y)=\lambda_2(y)\). Da
    \(\lambda_1,\lambda_2\) mit \(\lambda_1|_Z=\lambda_2|_Z\) beliebig waren,
    gilt \(\{x\}\bot_{\mathrm{st}}\{y\}\mid Z\).
\end{proof}

Wegen \(|L_{\mathrm{st}}|=2\) lässt sich die Rechteckbedingung weiter
konkretisieren. Wir betrachten die beiden Diagonalen
\[
    D_1=\{(\mathit{in},\mathit{in}),(\mathit{out},\mathit{out})\},
    \qquad
    D_2=\{(\mathit{in},\mathit{out}),(\mathit{out},\mathit{in})\}.
\]

\begin{Lemma}
    \label{lem:singl-diagonale}
    Eine Menge \(R\subseteq L_{\mathrm{st}}^2\) ist genau dann \emph{kein}
    Produktrechteck, wenn eine Diagonale
    \(D=\{(\alpha_1,\beta_1),(\alpha_2,\beta_2)\}\in\{D_1,D_2\}\) mit
    \(D\subseteq R\) existiert, der ein Kreuzpaar fehlt, d.\,h. mit
    \((\alpha_1,\beta_2)\notin R\) oder \((\alpha_2,\beta_1)\notin R\).
\end{Lemma}

\begin{proof}
    \label{pr:singl-diagonale}
    Enthält \(R\) eine Diagonale \(D\) ohne ein Kreuzpaar, so verletzt \(R\)
    die Produktbedingung und ist kein Rechteck. Ist umgekehrt \(R\) kein
    Produktrechteck, so existieren
    \((\alpha_1,\beta_1),(\alpha_2,\beta_2)\in R\) mit
    \((\alpha_1,\beta_2)\notin R\). Aus
    \((\alpha_1,\beta_1)\neq(\alpha_1,\beta_2)\) folgt \(\beta_1\neq\beta_2\),
    aus \((\alpha_1,\beta_2)\neq(\alpha_2,\beta_2)\) folgt
    \(\alpha_1\neq\alpha_2\). Wegen \(|L_{\mathrm{st}}|=2\) bilden
    \((\alpha_1,\beta_1),(\alpha_2,\beta_2)\) eine der Diagonalen
    \(D_1,D_2\), und \((\alpha_1,\beta_2)\) ist ein fehlendes Kreuzpaar.
\end{proof}

Proposition~\ref{prop:singl-rechteck} und Lemma~\ref{lem:singl-diagonale}
zeigen zusammen: Ein Gegenbeispiel zur Unabhängigkeit besteht genau aus einer
realisierbaren Belegung \(\eta\) und einer Diagonalen \(D\subseteq R_\eta\),
der ein Kreuzpaar fehlt. Dies rechtfertigt die Beschränkung auf
Kandidatenfälle mit \(\alpha_1\neq\alpha_2\) und \(\beta_1\neq\beta_2\): nur
solche Fälle können eine Rechteckverletzung bezeugen. Stimmen die Labels von
\(x\) (bzw. \(y\)) in beiden Labellings überein, so ist \(\lambda_3:=\lambda_2\)
(bzw. \(\lambda_3:=\lambda_1\)) bereits ein Zeuge der Unabhängigkeit.

\begin{Corollary}
    \label{cor:singl-fixlabel}
    Nimmt \(x\) oder \(y\) unter allen \(\lambda\in\Lambda_{\mathrm{st}}(F)\)
    nur ein einziges Label an, so gilt \(\{x\}\bot_{\mathrm{st}}\{y\}\mid Z\)
    für jedes \(Z\subseteq A\setminus\{x,y\}\).
\end{Corollary}

\begin{proof}
    \label{pr:singl-fixlabel}
    Nimmt \(x\) nur das Label \(\alpha\) an, so ist für jede realisierbare
    \(\eta\) jedes Paar in \(R_\eta\) von der Gestalt \((\alpha,\cdot)\). Damit
    ist \(R_\eta=\{\alpha\}\times Q\) mit
    \(Q=\{\,b\in L_{\mathrm{st}}\mid(\alpha,b)\in R_\eta\,\}\) ein
    Produktrechteck, und die Behauptung folgt mit
    Proposition~\ref{prop:singl-rechteck}. Der Fall für \(y\) verläuft
    symmetrisch.
\end{proof}

\section{CEGAR-style Algorithmus}
\label{cegar-alg}

\subsection{Realisierbarkeitsabfragen}
\label{subsec:singl-realize}

Der Algorithmus stützt sich auf eine einzige Hilfsabfrage, die prüft, ob sich
eine vorgegebene Teilbelegung zu einem \textit{stable}-Labelling fortsetzen
lässt; alle übrigen Schritte sind Schleifen über kombinatorische Belegungen.

Für ein Argument \(a\in A\) und ein Label \(\ell\in L_{\mathrm{st}}\) schreiben
wir \(\mathrm{Fix}_a(\ell)\) für das Literal, das \(a\) auf \(\ell\) festlegt:
\[
    \mathrm{Fix}_a(\mathit{in}):=i_a,
    \qquad
    \mathrm{Fix}_a(\mathit{out}):=\neg i_a,
\]
wobei \(i_a\) die \textit{in}-Variable aus Definition~\ref{def:st-pl} ist. Für
eine Teilbelegung \(\pi:W\to L_{\mathrm{st}}\) mit \(W\subseteq A\) definieren
wir die \emph{Realisierbarkeitsabfrage}
\[
    \mathrm{realisierbar}(\pi)
    \quad:\Longleftrightarrow\quad
    \mathrm{st}_F(\mathcal L)\land
    \bigwedge_{a\in W}\mathrm{Fix}_a(\pi(a))
    \text{ ist erfüllbar,}
\]
wobei \(\mathcal L=\{i_a\mid a\in A\}\) eine Kopie der Labelling-Variablen ist.
Nach Lemma~\ref{lem:st-pl-correct} gilt \(\mathrm{realisierbar}(\pi)\) genau
dann, wenn ein \(\lambda\in\Lambda_{\mathrm{st}}(F)\) mit \(\lambda|_W=\pi\)
existiert. Praktisch ist dies ein einzelner SAT-Aufruf auf
\(\mathrm{st}_F(\mathcal L)\) mit den Assumptions
\(\{\,i_a\mid \pi(a)=\mathit{in}\,\}\cup
\{\,\neg i_a\mid \pi(a)=\mathit{out}\,\}\); einzig die SAT/UNSAT-Antwort wird
benötigt. Insbesondere ist \(\mathrm{realisierbar}(\emptyset)\) genau dann
wahr, wenn \(\Lambda_{\mathrm{st}}(F)\neq\emptyset\).

Für die Enumeration gemeinsamer \(Z\)-Belegungen verwenden wir zusätzlich zwei
Kopien \(\mathcal L^1,\mathcal L^2\) der Labelling-Variablen. Wir schreiben
\(\mathrm{Fix}^k_a(\ell)\) für das entsprechende Literal über \(\mathcal L^k\)
(also \(i_a^k\) bzw. \(\neg i_a^k\)) und verwenden die Gleichheitsformel
\(\mathrm{Eq}^{1,2}_Z\) aus Definition~\ref{def:eq-qbf}. Da alle Abfragen
dieselbe Kodierung \(\mathrm{st}_F\) teilen und sich nur in Assumptions und
Blocking-Klauseln unterscheiden, lassen sie sich inkrementell beantworten.

\subsection{Algorithmus}
\label{subsec:singl-algo}

Nach Proposition~\ref{prop:singl-rechteck} ist zu prüfen, ob \(R_\eta\) für
jede realisierbare Belegung \(\eta\) ein Produktrechteck ist. Nach
Lemma~\ref{lem:singl-diagonale} genügt es, die beiden Diagonalen
\(D_1,D_2\) zu betrachten und für jede gemeinsam realisierbare Diagonale ein
fehlendes Kreuzpaar zu suchen.

\begin{enumerate}
    \item \textbf{Globale Erfüllbarkeit.} Ist
    \(\mathrm{realisierbar}(\emptyset)\) falsch, so gilt
    \(\Lambda_{\mathrm{st}}(F)=\emptyset\). Damit ist
    \(\{x\}\bot_{\mathrm{st}}\{y\}\mid Z\) vakuos erfüllt; der Algorithmus
    antwortet, dass Unabhängigkeit gilt.

    \item \textbf{Fixed-Label-Vortest.} Prüfe für \(a\in\{x,y\}\) und
    \(\ell\in L_{\mathrm{st}}\) jeweils \(\mathrm{realisierbar}(a\mapsto\ell)\).
    Nimmt \(x\) oder \(y\) nur ein einziges Label an, so gilt nach
    Korollar~\ref{cor:singl-fixlabel} Unabhängigkeit. Dieser Test kostet
    höchstens vier SAT-Aufrufe und entscheidet ganze Instanzklassen vorab.

    \item \textbf{Diagonalenumeration.} Für jede Diagonale
    \(D=\{(\alpha_1,\beta_1),(\alpha_2,\beta_2)\}\in\{D_1,D_2\}\):
    \begin{enumerate}
        \item Setze \(B:=\top\) und bilde
        \[
        \begin{aligned}
            C^{12}_{D}
            :=\ &
            \mathrm{st}_F(\mathcal L^1)\land\mathrm{st}_F(\mathcal L^2)
            \land\mathrm{Fix}^1_x(\alpha_1)\land\mathrm{Fix}^1_y(\beta_1)\\
            &\land\mathrm{Fix}^2_x(\alpha_2)\land\mathrm{Fix}^2_y(\beta_2)
            \land\mathrm{Eq}^{1,2}_Z.
        \end{aligned}
        \]
        \item \textbf{Gemeinsame \(Z\)-Belegung suchen.} Solange
        \(C^{12}_{D}\land B\) erfüllbar ist, wähle ein Modell und lies die
        gemeinsame Belegung \(\eta=\lambda_1|_Z=\lambda_2|_Z\) ab. Nach
        Konstruktion gilt \(D\subseteq R_\eta\).
        \item \textbf{Kreuzpaare prüfen.} Ist
        \[
            \neg\,\mathrm{realisierbar}\bigl(\{x\mapsto\alpha_1,\,y\mapsto\beta_2\}\cup\eta\bigr)
            \quad\text{oder}\quad
            \neg\,\mathrm{realisierbar}\bigl(\{x\mapsto\alpha_2,\,y\mapsto\beta_1\}\cup\eta\bigr),
        \]
        so fehlt \(R_\eta\) ein Kreuzpaar; nach
        Lemma~\ref{lem:singl-diagonale} ist \(R_\eta\) kein Produktrechteck.
        Der Algorithmus gibt das Gegenbeispiel \((\eta,D)\) zurück und
        terminiert mit \(\{x\}\not\bot_{\mathrm{st}}\{y\}\mid Z\).
        \item \textbf{Blocking.} Sind beide Kreuzpaare realisierbar, so ist
        \(R_\eta=L_{\mathrm{st}}^2\) ein Produktrechteck. Schließe \(\eta\)
        durch die Blocking-Klausel
        \[
            \mathrm{Block}(\eta):=\bigvee_{z\in Z}\neg\,\mathrm{Fix}^1_z(\eta(z))
        \]
        aus, setze \(B:=B\land\mathrm{Block}(\eta)\) und fahre mit der
        nächsten gemeinsamen \(Z\)-Belegung fort.
    \end{enumerate}

    \item \textbf{Abschluss.} Wurden beide Diagonalen ohne Gegenbeispiel
    durchlaufen, so ist nach Proposition~\ref{prop:singl-rechteck} jede
    realisierbare Belegung ein Produktrechteck, und es gilt
    \(\{x\}\bot_{\mathrm{st}}\{y\}\mid Z\).
\end{enumerate}

Dabei verwenden wir die Konvention \(\bigvee\emptyset=\bot\). Im Fall
\(Z=\emptyset\) ist \(\mathrm{Block}(\eta)=\bot\), sodass jede Diagonale nach
genau einer Iteration abbricht.

\subsection{Korrektheit und Terminierung}
\label{subsec:singl-korrekt}

\begin{Lemma}
    \label{lem:singl-korrekt}
    Der Algorithmus aus Abschnitt~\ref{subsec:singl-algo} entscheidet
    \(\textit{SingInd}_{\mathrm{st}}\) korrekt und terminiert.
\end{Lemma}

\begin{proof}
    \label{pr:singl-korrekt}
    \emph{Korrektheit bei Gegenbeispiel.} Gibt der Algorithmus in
    Schritt~3(c) ein Gegenbeispiel \((\eta,D)\) zurück, so wurde \(\eta\) als
    gemeinsame \(Z\)-Belegung mit \(D\subseteq R_\eta\) gefunden, und ein
    Kreuzpaar fehlt in \(R_\eta\). Nach Lemma~\ref{lem:singl-diagonale} ist
    \(R_\eta\) kein Produktrechteck; mit Proposition~\ref{prop:singl-rechteck}
    folgt \(\{x\}\not\bot_{\mathrm{st}}\{y\}\mid Z\).

    \emph{Korrektheit bei Unabhängigkeit.} Meldet der Algorithmus
    Unabhängigkeit, so liegt einer von drei Fällen vor:
    \(\Lambda_{\mathrm{st}}(F)=\emptyset\) (Schritt~1, vakuos erfüllt), ein
    fixiertes Label nach Korollar~\ref{cor:singl-fixlabel} (Schritt~2) oder
    der erschöpfende Durchlauf beider Diagonalen (Schritt~4). Im letzten Fall
    enthält kein \(R_\eta\) eine Diagonale ohne Kreuzpaar; nach
    Lemma~\ref{lem:singl-diagonale} ist jedes \(R_\eta\) ein Produktrechteck,
    und Proposition~\ref{prop:singl-rechteck} liefert
    \(\{x\}\bot_{\mathrm{st}}\{y\}\mid Z\).

    \emph{Vollständigkeit.} Gilt \(\{x\}\not\bot_{\mathrm{st}}\{y\}\mid Z\), so
    existiert nach Proposition~\ref{prop:singl-rechteck} eine realisierbare
    Belegung \(\eta^\ast\), für die \(R_{\eta^\ast}\) kein Produktrechteck ist.
    Nach Lemma~\ref{lem:singl-diagonale} gibt es eine Diagonale
    \(D\subseteq R_{\eta^\ast}\) mit fehlendem Kreuzpaar. Im Durchlauf dieser
    Diagonalen erfüllt \(\eta^\ast\) die Formel \(C^{12}_{D}\); geblockt werden
    ausschließlich Belegungen mit \(R_\eta=L_{\mathrm{st}}^2\), also nicht
    \(\eta^\ast\). Folglich erreicht die Enumeration \(\eta^\ast\) (oder zuvor
    ein anderes Gegenbeispiel), und die Kreuzpaarprüfung schlägt fehl.

    \emph{Terminierung.} In jeder Iteration der Enumerationsschleife wird
    entweder terminiert oder genau eine realisierbare gemeinsame
    \(Z\)-Belegung durch \(\mathrm{Block}(\eta)\) dauerhaft ausgeschlossen.
    Davon gibt es höchstens \(2^{|Z|}\); die äußere Schleife über die beiden
    Diagonalen ist endlich.
\end{proof}

\begin{Corollary}
    \label{cor:singl-aufrufe}
    Sei \(N\) die Anzahl der in Schritt~3 insgesamt aufgezählten gemeinsamen
    \(Z\)-Belegungen. Der Algorithmus benötigt \(N+O(1)\) Enumerationsabfragen
    und höchstens \(2N+O(1)\) weitere Realisierbarkeitsabfragen. Im
    schlechtesten Fall ist \(N\in O(2^{|Z|})\).
\end{Corollary}

\begin{proof}
    \label{pr:singl-aufrufe}
    Globale Erfüllbarkeit und Fixed-Label-Vortest verursachen höchstens fünf
    Aufrufe. Jede der \(N\) aufgezählten Belegungen kostet eine
    \(C^{12}\)-Abfrage und höchstens zwei Kreuzpaarabfragen; pro Diagonale
    kommt eine abschließende UNSAT-Abfrage hinzu. Die Schranke
    \(N\in O(2^{|Z|})\) folgt aus der Zahl möglicher \(Z\)-Belegungen.
\end{proof}

Die exponentielle Schranke ist der \(\Pi_2^{\mathrm P}\)-Härte des Problems
geschuldet (Theorem~\ref{th:ind-st-is-pi2p-c}). Die folgenden Heuristiken
zielen darauf, die Zahl \(N\) der aufgezählten Belegungen sowie die Kosten der
Einzelaufrufe zu senken.

\subsection{Heuristiken zur Gegenbeispielsuche}
\label{subsec:singl-heuristiken}

Die folgenden Heuristiken senken Anzahl bzw. Kosten der SAT-Aufrufe, ohne die
Korrektheit zu beeinflussen.

\paragraph{Inkrementelles SAT mit Assumptions.}
Sämtliche Anfragen teilen die Basisformel \(\mathrm{st}_F\) (bzw. ihre beiden
Kopien) und unterscheiden sich nur in Assumptions und Blocking-Klauseln. Ein
einziger persistenter Solver über alle Diagonalen und Belegungen hinweg
erlaubt die Wiederverwendung gelernter Klauseln und vermeidet wiederholtes
Aufbauen der Formel.

\paragraph{Blocking über UNSAT-Cores.}
Ist eine Kreuzpaarabfrage unter den Assumptions einer Belegung \(\eta\)
unerfüllbar, liefert der Solver einen UNSAT-Core, also eine Teilmenge der
\(\eta\)-Literale, die den Widerspruch bereits erzwingt. Dessen Negation
blockt nicht nur \(\eta\), sondern alle Belegungen mit derselben
konfliktären Teilbelegung und beschneidet den Suchraum oft deutlich stärker
als das punktweise \(\mathrm{Block}(\eta)\).

\paragraph{Verallgemeinertes Blocking.}
Statt nur \(\eta\) auszuschließen, kann aus dem Modell einer erfüllbaren
Kreuzpaarprüfung eine ganze Familie unkritischer Belegungen abgeleitet werden,
etwa durch Weglassen von Literalen, die für die Realisierbarkeit der
Kreuzpaare irrelevant sind. Dies entspricht der Verfeinerungsphase eines
echten CEGAR-Verfahrens für die \(\exists\forall\)-Struktur des Problems.

\paragraph{Globales Backbone von \(Z\).}
Argumente \(z\in Z\), die über alle \textit{stable}-Labellings hinweg dasselbe
Label tragen (das \emph{Backbone} von \(\mathrm{st}_F\)), sind in jeder
gemeinsamen \(Z\)-Belegung fixiert. Wird dieses Backbone einmal vorab
berechnet, lässt sich die Enumeration auf die variablen Argumente
\(Z_{\mathrm{var}}\subseteq Z\) beschränken und die Blocking-Klauseln
verkürzen. Im Gegensatz zu einer kandidatenweisen Determinismusanalyse fällt
dieser Aufwand nur einmal an.

\paragraph{SCC-Vorverarbeitung.}
Zerfällt \(F\) in mehrere starke Zusammenhangskomponenten, propagieren
\textit{stable}-Labels entlang des SCC-DAG. Argumente außerhalb der für \(x\),
\(y\) und \(Z\) relevanten Komponenten sind häufig bereits fixiert, wodurch
sich die Formeln verkleinern. Bildet das relevante Fragment einen einzelnen
geraden Zykel, ist die Frage ohne SAT-Aufruf in linearer Zeit entscheidbar.

\paragraph{Reihenfolge der Abfragen.}
Innerhalb eines Diagonaldurchlaufs sollte zuerst das voraussichtlich
diskriminierende Kreuzpaar geprüft werden, um ein Gegenbeispiel früh zu
finden. Durch \emph{phase saving} im Solver lassen sich strukturell
unterschiedliche \(Z\)-Belegungen schneller erreichen, was \(N\) in der Praxis
senkt.